% Included by literature_review.tex; edit definitions here.
\clearpage
\pdfbookmark[0]{Medical and neuroscience nomenclature}{medical-nomenclature}
\phantomsection\label{sec:medical-nomenclature}
\section*{Medical and neuroscience nomenclature}
Plain-English explanations of the medical and neuroscience terms used
in this review, the paper notes and the related reading. Entries are
grouped by topic. Frequency bands are approximate: individual papers
may use different boundaries.

\begingroup
\small
\setlength{\LTpre}{5pt}
\setlength{\LTpost}{5pt}
\newcommand{\medicalterm}[2]{\textbf{#1} & #2\\[4pt]}
\newcommand{\medicalgroup}[1]{\multicolumn{2}{@{}l}{\rule{0pt}{2.2em}\large\bfseries #1}\\*[6pt]}
\begin{longtable}{@{}>{\raggedright\arraybackslash}p{0.29\linewidth}>{\raggedright\arraybackslash}p{\dimexpr0.71\linewidth-2\tabcolsep\relax}@{}}
\textbf{Term / abbreviation} & \textbf{Meaning}\\\hline
\endfirsthead
\textbf{Term / abbreviation} & \textbf{Meaning (continued)}\\\hline
\endhead
\endfoot
\endlastfoot

\medicalgroup{Neurons and cellular activity}
\medicalterm{Action potential / spike}{A brief electrical pulse generated by a neuron and carried along its axon.}
\medicalterm{Axon}{The extension of a neuron that carries spikes towards other cells.}
\medicalterm{Burst / bursting}{A cluster of spikes occurring close together in time.}
\medicalterm{Cell body / soma}{The main part of a neuron, containing its cell nucleus.}
\medicalterm{Dendrite}{A branching extension through which a neuron receives many of its inputs.}
\medicalterm{Depolarization}{A change that makes the voltage inside a neuron less negative relative to its surroundings.}
\medicalterm{Excitation / excitatory}{An effect that makes a neuron more likely to fire. An excitatory neuron usually has this effect on its targets.}
\medicalterm{Fast-spiking neuron / FSI}{A neuron capable of sustained rapid firing. FSI means fast-spiking interneuron, commonly an inhibitory cell.}
\medicalterm{Firing}{A neuron producing action potentials.}
\medicalterm{Firing rate}{The number of spikes per unit time, usually expressed as spikes per second. It is estimated over a time window.}
\medicalterm{Hyperpolarization}{A change that makes the voltage inside a neuron more negative relative to its surroundings.}
\medicalterm{Inhibition / inhibitory}{An effect that makes a neuron less likely to fire; it need not stop firing completely.}
\medicalterm{Interneuron}{A neuron whose connections mainly influence a local circuit. Many cortical interneurons are inhibitory.}
\medicalterm{Membrane potential}{The voltage difference between the inside and outside of a cell.}
\medicalterm{Neural dynamics}{How neural activity changes over time through inputs and interactions within a circuit.}
\medicalterm{Neural circuit / network}{Interconnected neurons or populations that influence one another.}
\medicalterm{Neural population / ensemble}{A group of neurons considered together, for example because they share a location, cell type or activity pattern.}
\medicalterm{Neuron}{A nerve cell that receives, processes and transmits signals.}
\medicalterm{Neurophysiology}{The study of how neurons and the nervous system function.}
\medicalterm{Parvalbumin neuron / PV neuron}{A neuron containing the calcium-binding protein parvalbumin. Cortical PV neurons are commonly fast-spiking inhibitory interneurons.}
\medicalterm{Population firing rate}{The average firing rate of a group of neurons. A recorded sample can estimate it if the sample represents the group.}
\medicalterm{Pyramidal neuron}{A common excitatory cortical neuron named after its roughly pyramid-shaped cell body; ``pyramidal'' names the cell type, not an oscillation.}
\medicalterm{Spontaneous / endogenous / native activity}{Activity generated by the neural system itself, rather than directly prescribed by an external stimulus.}
\medicalterm{Spike train}{The sequence of times at which a neuron fires.}
\medicalterm{Tonic / phasic activity}{Tonic activity is relatively sustained; phasic activity consists of brief changes or responses.}

\medicalgroup{Synapses, signalling and adaptation}
\medicalterm{Dopamine}{A chemical messenger that changes how neural circuits operate, including circuits involved in movement, motivation and learning.}
\medicalterm{Dopaminergic}{Producing, releasing, responding to or otherwise involving dopamine.}
\medicalterm{Dopamine depletion / hypodopaminergic}{Reduced dopamine availability or signalling in the relevant circuit.}
\medicalterm{D1 / D2 receptors}{Two families of proteins through which dopamine affects cells. Their effects depend on the cell and circuit.}
\medicalterm{GABA}{Gamma-aminobutyric acid: the main inhibitory chemical messenger in the adult brain.}
\medicalterm{GABAergic}{Using GABA to transmit signals.}
\medicalterm{Glutamate / glutamatergic}{Glutamate is the main excitatory chemical messenger in the brain; glutamatergic means using glutamate.}
\medicalterm{Neurotransmitter}{A chemical released by a neuron to communicate with another cell.}
\medicalterm{Postsynaptic potential}{A change in a receiving neuron's membrane voltage caused by synaptic input. It differs from a spike.}
\medicalterm{Presynaptic / postsynaptic}{Presynaptic refers to the sending side of a synapse; postsynaptic refers to the receiving side.}
\medicalterm{Receptor}{A protein that responds to a signalling molecule and changes the receiving cell's activity.}
\medicalterm{Synapse / synaptic}{A connection through which one neuron influences another cell; synaptic means relating to that connection.}
\medicalterm{Synaptic depression}{A decrease in the effect of a synapse during repeated activity. Here ``depression'' does not mean a mood disorder.}
\medicalterm{Synaptic plasticity}{A change in the strength or effectiveness of a synapse over time.}

\medicalgroup{Brain regions and pathways}
\medicalterm{Afferent / efferent}{Afferent connections arrive at a specified region; efferent connections leave it.}
\medicalterm{Arkypallidal neuron / Arky}{A type of GPe neuron with prominent projections back to the striatum.}
\medicalterm{Basal ganglia / BG}{Interconnected brain structures involved in selecting and regulating actions, as well as other functions.}
\medicalterm{Cortex / cortical}{The brain's outer sheet of neural tissue; cortical means relating to it.}
\medicalterm{Cortico-subthalamic}{Connecting or involving the cortex and the subthalamic nucleus.}
\medicalterm{Contralateral / ipsilateral}{On the opposite side / the same side of the body relative to the structure or intervention being described.}
\medicalterm{D1-MSN / D2-MSN}{Striatal neurons distinguished by dopamine-receptor expression, typically associated with the direct and indirect pathways respectively.}
\medicalterm{Direct pathway}{A basal-ganglia route from the striatum to output nuclei such as GPi/SNr, without passing through GPe.}
\medicalterm{Dorsal / ventral}{In the spinal cord, towards the back / the belly side. Dorsoventral means spanning this direction.}
\medicalterm{Fibre / tract / projection}{An axon, bundle of axons or connection carrying signals between neurons or regions.}
\medicalterm{GPe}{External globus pallidus: a basal-ganglia region interacting with the STN, striatum and other regions.}
\medicalterm{GPi}{Internal globus pallidus: a major basal-ganglia output region and a DBS target. It differs from GPe.}
\medicalterm{Hippocampus}{A brain structure important for memory and spatial navigation.}
\medicalterm{Hyperdirect pathway}{A route carrying signals directly from cortex to STN, bypassing the striatum.}
\medicalterm{Indirect pathway}{A basal-ganglia route involving striatum, GPe and downstream output circuits; the classical route also passes through STN.}
\medicalterm{Motor cortex / M1}{Cortical regions involved in movement; M1 is the primary motor cortex.}
\medicalterm{MSN / SPN}{Medium spiny neuron / spiny projection neuron: a principal inhibitory projection neuron of the striatum.}
\medicalterm{Nucleus}{In brain anatomy, a group of neuron cell bodies. This differs from the nucleus inside an individual cell.}
\medicalterm{Pallidal / pallidostriatal}{Pallidal means relating to the globus pallidus; pallidostriatal refers to connections from pallidum to striatum.}
\medicalterm{Premotor cortex / SMA}{Cortical areas involved in preparing and organizing movements; SMA is the supplementary motor area.}
\medicalterm{Prototypic neuron / Proto}{A major GPe cell type that projects to other basal-ganglia regions, including STN.}
\medicalterm{Rostral / caudal}{Towards the head / towards the tail end of the body or spinal cord.}
\medicalterm{Sensorimotor}{Relating to sensory information and movement control.}
\medicalterm{SNc}{Substantia nigra pars compacta: a region containing dopamine-producing neurons that are particularly affected in Parkinson's disease.}
\medicalterm{SNr}{Substantia nigra pars reticulata: a major basal-ganglia output region, distinct from SNc.}
\medicalterm{Spinal cord}{Neural tissue connecting the brain with the body and containing circuits involved in movement and sensation.}
\medicalterm{STN}{Subthalamic nucleus: a small, mainly excitatory basal-ganglia structure and a common DBS target.}
\medicalterm{Striatum / striatal}{A major input region of the basal ganglia; striatal means relating to it.}
\medicalterm{Subcortical}{Located beneath the cerebral cortex.}
\medicalterm{Thalamus / thalamocortical}{The thalamus is a group of brain nuclei communicating with cortex and other regions; thalamocortical means connecting thalamus and cortex.}
\medicalterm{Vim}{Ventral intermediate nucleus of the thalamus: a stimulation target used in tremor treatment.}

\medicalgroup{Neural signals and rhythms}
\medicalterm{Alpha / theta activity}{Named rhythm bands: alpha is roughly 8--12 Hz; theta is roughly 4--8 Hz, extended to 10 Hz in some reviewed work.}
\medicalterm{Beta activity}{Rhythmic activity at roughly 13--30 Hz; some reviewed papers extend the band to 35 Hz.}
\medicalterm{Coherence}{A frequency-specific measure of how consistently two signals are related. It does not by itself establish causation.}
\medicalterm{Entrainment}{A rhythm becoming aligned with an external rhythm, such as periodic stimulation.}
\medicalterm{Evoked response}{Neural activity occurring in response to a stimulus.}
\medicalterm{Gamma activity}{Rhythmic activity above roughly 30 Hz. Upper limits and the definition of high gamma vary by study.}
\medicalterm{Granger causality}{A test of whether one signal's past improves prediction of another. The name does not imply proof of a biological cause.}
\medicalterm{LFP}{Local field potential: a voltage recorded in neural tissue, reflecting combined nearby electrical currents, especially synaptic activity.}
\medicalterm{Neural oscillation / rhythm}{A recurring rise and fall in neural activity. It can occur in population rates or field potentials.}
\medicalterm{PAC}{Phase-amplitude coupling: a statistical relationship between the phase of a slower rhythm and the amplitude of faster activity.}
\medicalterm{Phase / amplitude}{Phase is position within an oscillation's cycle; amplitude describes the size of its variation.}
\medicalterm{Readiness potential}{A slow electrical signal associated with preparation for voluntary movement.}
\medicalterm{Spectral / band power}{How much signal variation lies within a frequency range. Beta power, for example, measures variation in the beta band.}
\medicalterm{Synchronization / synchrony}{Coordination in the timing of neural activity. Neurons need not have identical average firing rates.}
\medicalterm{Wilson--Cowan / firing-rate model}{A population-level description of how neural activity changes through excitation, inhibition and external inputs; Wilson--Cowan models commonly use coupled excitatory and inhibitory populations.}

\medicalgroup{Recording and imaging methods}
\medicalterm{Calcium imaging / indicator}{Recording activity-related changes in intracellular calcium using a fluorescent sensor; calcium signals are slower than individual spikes.}
\medicalterm{Connectome / connectivity}{A map or description of connections between brain regions. Structural connectivity concerns anatomy; functional connectivity concerns statistical relationships between activity signals.}
\medicalterm{ECoG}{Electrocorticography: recording electrical brain activity using electrodes on the cortical surface.}
\medicalterm{EEG}{Electroencephalography: recording electrical brain activity, usually with electrodes on the scalp.}
\medicalterm{Electrode / microelectrode}{A conducting contact used to record or stimulate electrically; a microelectrode is a small contact suited to localized recordings.}
\medicalterm{Electrophysiology}{Measurement of electrical activity in cells or tissue.}
\medicalterm{EMG}{Electromyography: recording electrical activity associated with muscle activation.}
\medicalterm{Extracellular recording}{Recording from outside cells; depending on the setup, this can capture spikes and field potentials.}
\medicalterm{Intracranial recording}{Recording using electrodes placed inside the skull.}
\medicalterm{MEG}{Magnetoencephalography: measuring magnetic fields produced by brain activity.}
\medicalterm{MRI / fMRI}{Magnetic resonance imaging measures tissue properties; functional MRI uses blood-oxygen-related changes to study brain activity indirectly.}
\medicalterm{Multi-unit / single-unit activity}{Multi-unit activity combines spikes from several neurons; single-unit activity consists of spikes assigned to one neuron.}
\medicalterm{PET}{Positron emission tomography: imaging the distribution of an injected radioactive tracer, for example to study molecular systems.}
\medicalterm{Polysomnography / PSG}{A sleep recording combining signals such as EEG, eye movements, muscle activity and breathing.}
\medicalterm{sEEG}{Stereo-electroencephalography: recording electrical activity with depth electrodes placed within the brain.}
\medicalterm{Spike sorting}{Assigning recorded spikes to putative individual neurons using their waveform and other features.}
\medicalterm{Tractography}{Estimating the courses of nerve-fibre bundles from diffusion MRI data.}

\medicalgroup{Stimulation and optogenetics}
\medicalterm{Adaptive DBS / aDBS}{DBS whose settings change automatically in response to measured signals.}
\medicalterm{Antidromic / orthodromic}{Antidromic spikes travel towards the cell body along an axon; orthodromic spikes travel in its usual output direction.}
\medicalterm{Archaerhodopsin / Arch / ArchT}{Light-driven proton pumps used to inhibit neurons; ArchT is an engineered variant.}
\medicalterm{Channelrhodopsin / ChR2 / Chronos}{Light-sensitive ion channels used to excite neurons; ChR2 and Chronos are examples.}
\medicalterm{Closed-loop stimulation}{Stimulation chosen using ongoing measurements. Its purpose may be to increase, suppress or otherwise study activity.}
\medicalterm{DBS}{Deep brain stimulation: electrical stimulation delivered through implanted electrodes to influence brain circuits.}
\medicalterm{Halorhodopsin / NpHR / eNpHR}{A light-driven chloride pump used to inhibit neurons; eNpHR denotes enhanced versions.}
\medicalterm{HFS}{High-frequency stimulation: a rapid train of stimulation pulses; the exact frequency depends on the experiment.}
\medicalterm{Neuromodulation}{Changing neural activity using an intervention such as electrical stimulation, optical stimulation or a chemical agent.}
\medicalterm{Open-loop stimulation}{Stimulation delivered without adjusting it in response to ongoing measurements of the controlled activity.}
\medicalterm{Opsin}{A light-sensitive protein. In these experiments, introduced opsins allow light to influence neuronal activity.}
\medicalterm{Optogenetics}{Using genetically introduced light-sensitive proteins to manipulate selected cells or their connections with light.}
\medicalterm{Optrode}{A device combining optical stimulation with electrode recording.}
\medicalterm{Photostimulation}{Stimulation using light.}
\medicalterm{Projection targeting}{Selecting neurons or their axons according to which region they connect to.}
\medicalterm{Pulse width / stimulation amplitude}{Pulse width is a pulse's duration; amplitude is its electrical or optical strength.}
\medicalterm{Stimulation artifact}{A recorded signal caused by the stimulation or equipment rather than the neural activity being studied.}
\medicalterm{Viral vector / opsin expression}{A modified virus can deliver genetic instructions; opsin expression means cells produce the instructed light-sensitive protein.}

\medicalgroup{Symptoms, disorders and treatment terms}
\medicalterm{Akinesia}{Difficulty initiating movement, or a marked lack of movement.}
\medicalterm{Axial symptoms}{Problems mainly affecting the trunk, posture, balance or walking.}
\medicalterm{Biomarker}{A measurable biological characteristic used to indicate a condition or its response to an intervention.}
\medicalterm{Bradykinesia}{Slowness of movement, often with progressively smaller or slower repeated movements.}
\medicalterm{Cognitive}{Relating to processes such as attention, memory, learning and decision-making.}
\medicalterm{Dyskinesia}{Involuntary movements; in Parkinson's research, often abnormal movements associated with dopaminergic treatment.}
\medicalterm{Dystonia}{Involuntary muscle contractions causing abnormal postures or repetitive twisting movements.}
\medicalterm{Epilepsy / seizure}{A seizure is an episode of abnormal brain activity; epilepsy is a disorder involving an enduring tendency to have seizures.}
\medicalterm{Essential tremor}{A movement disorder characterized mainly by tremor during posture or action.}
\medicalterm{Freezing of gait / FoG}{A brief inability or marked difficulty to step despite the intention to walk.}
\medicalterm{Gait}{The pattern of walking.}
\medicalterm{Huntington's disease}{An inherited neurodegenerative disorder affecting movement, cognition and behaviour.}
\medicalterm{Hypokinetic / hyperkinetic}{Hypokinetic refers to reduced movement; hyperkinetic refers to excessive or unwanted movement.}
\medicalterm{Levodopa / L-DOPA}{A substance converted into dopamine in the body and used to treat Parkinson's symptoms.}
\medicalterm{Motor / non-motor symptoms}{Motor symptoms affect movement; non-motor symptoms include effects on sleep, mood, thinking and other functions.}
\medicalterm{Neurodegeneration}{Progressive damage to or loss of neurons.}
\medicalterm{ON / OFF medication}{Study conditions describing treatment benefit or medication exposure; the paper specifies which meaning is used. These differ from stimulation ON/OFF.}
\medicalterm{Parkinson's disease / PD}{A progressive neurodegenerative disorder involving dopamine-neuron loss and motor as well as non-motor symptoms.}
\medicalterm{Parkinsonism / parkinsonian}{Parkinsonism is a pattern of movement signs resembling Parkinson's disease; parkinsonian also describes these features in experimental models.}
\medicalterm{Pathological / physiological}{Pathological concerns disease-related dysfunction; physiological concerns normal biological function.}
\medicalterm{Phenotype}{An observable characteristic, such as a movement pattern or disease-related behaviour.}
\medicalterm{Rigidity}{Increased resistance when a joint is moved by someone else; commonly experienced as stiffness.}
\medicalterm{Tourette syndrome}{A condition involving recurrent motor and vocal tics.}
\medicalterm{Tremor}{An involuntary rhythmic shaking movement.}
\medicalterm{Volitional / voluntary movement}{Movement initiated intentionally.}

\medicalgroup{Sleep and experimental terminology}
\medicalterm{Acute / chronic}{Acute usually refers to a short-term condition or experiment; chronic refers to a persistent condition or long-term recording/intervention.}
\medicalterm{Anaesthetized}{Under anaesthesia; the drugs can affect neural activity and responses to stimulation.}
\medicalterm{Brain slice}{A thin piece of brain tissue kept viable outside the animal for experiments.}
\medicalterm{Biophysical}{Describing biological activity through physical mechanisms such as membrane currents and ion channels.}
\medicalterm{Causal / correlational evidence}{Causal evidence supports one process influencing another; correlation describes variables changing together.}
\medicalterm{Histology}{Examining tissue structure, usually with microscopy and stains.}
\medicalterm{Epidural}{Located outside the dura mater, the tough outer membrane surrounding the brain and spinal cord.}
\medicalterm{In silico / in vitro / in vivo}{In silico: in a computer model. In vitro: outside the living organism. In vivo: within a living organism.}
\medicalterm{NHP}{Non-human primate, such as a monkey used in research.}
\medicalterm{NREM sleep / N1--N3}{Non-rapid-eye-movement sleep, subdivided into stages N1, N2 and N3.}
\medicalterm{REM sleep}{Rapid-eye-movement sleep: a sleep state with characteristic brain activity, eye movements and usually reduced muscle tone.}
\medicalterm{Sleep decoding}{Estimating sleep state from recorded signals using an algorithm.}
\medicalterm{Transgenic animal}{An animal whose genetic material has been deliberately altered, for example to express an opsin in selected neurons.}
\medicalterm{6-OHDA / MPTP model}{Experimental Parkinsonian models using substances that damage dopamine systems; their effects do not reproduce every feature of human disease.}
\end{longtable}

\textbf{Sources.} Circuit, recording and optogenetic terminology follows
the reviewed papers, particularly
\cite{grosenick2015,fenno2011,neumann2023,delacrompe2025}.
Additional plain-language checks:
\href{https://www.ninds.nih.gov/health-information/disorders/glossary-neurological-terms}{NINDS neurological glossary},
\href{https://www.ninds.nih.gov/current-research/focus-disorders/parkinsons-disease-research/parkinsons-disease-challenges-progress-and-promise}{NINDS Parkinson's disease overview}, and
\href{https://www.ncbi.nlm.nih.gov/books/NBK11117/}{\emph{Neuroscience}: excitatory and inhibitory postsynaptic potentials}.
\endgroup
